\documentclass[10pt]{article}

% Specify the margins and text size.
\setlength{\textwidth}{6.4in}
\setlength{\textheight}{9.5in}
\setlength{\oddsidemargin}{0pt}
\setlength{\evensidemargin}{0pt}
\setlength{\topmargin}{0pt}
\setlength{\hoffset}{.05in}
\setlength{\voffset}{-.8in}

\setlength{\parskip}{5pt}
\setlength{\parindent}{0pt}

% Load some fonts and symbol packages
\usepackage{latexsym}
\usepackage{pifont}       % contains 'star' symbol for counterinsurgency handbook title
\usepackage{yfonts} 
\usepackage{amsmath}
\usepackage{amsfonts}

\usepackage{graphicx}     % actually, this is loaded by pstricks
\usepackage[T1]{fontenc}
\usepackage{ifthen}
\usepackage{pstricks,pst-func,pst-text,pst-node,multido,pst-plot,calc,pst-3dplot}
%\usepackage[all]{xy}
%\usepackage{animate}

% The hyperref package inserts links into the pdf.
\definecolor{MyLinkColor}{rgb}{.1,.2,1}
\definecolor{MyCiteColor}{rgb}{.1,1,.2}
\definecolor{MyURLColor}{rgb}{.4,.4,.4}
\usepackage[backref=true,pagebackref=false,hyperindex,colorlinks=true,
  linkcolor=MyLinkColor,urlcolor=MyURLColor]{hyperref}


% The tweaklist package is something I found on the web.  It provides a simple interface
% for making changes to spacing used in the itemize and enumerate environments.  Comment
% this out if you don't care to use tweaklist.
\usepackage{tweaklist}
\renewcommand{\itemhook}{\setlength{\parskip}{2pt}\setlength{\parsep}%
{1pt}\setlength{\topsep}{0pt}\setlength{\itemsep}{0pt}}

\newcommand{\N}{\ensuremath\mathbb{N}}
\newcommand{\Q}{\ensuremath\mathbb{Q}}
\newcommand{\R}{\ensuremath\mathbb{R}}
\newcommand{\C}{\ensuremath\mathbb{C}}
\newcommand{\z}{\hphantom{0}}

\begin{document}
\pagestyle{empty}

\href{http://www.su.edu}{\includegraphics[height=1.75cm]{sulogo.eps}}
\vspace{-1.69cm}

{\small\hfill
\begin{tabular}{cl}
 & Math 207\\ &  Introduction to Statistics\\
\end{tabular}

\begin{center}
\textbf{\large Population and Sample Standard Deviations}
\end{center}

If $x_1$, $x_2$, $\dots$, $x_n$ are the data values for an entire population, 
the \textbf{population standard deviation} is 
\[\mbox{SD} = \sqrt{\frac{1}{n}\,\sum_{i=1}^n\,\left(x_i - \mbox{mean}_x\right)^2}\]
If the data values are only a sample from the population, then the
\textbf{sample standard deviation} is
\[\mbox{sd} = \sqrt{\frac{1}{n-1}\,\sum_{i=1}^n\,\left(x_i - \mbox{mean}_x\right)^2}\]
\bigskip

1. Explain in words what the difference is in the two formulas.  What happens if the sample
size $n$ is large?
\vspace{1in}

2. Each of the following lists has an average of 50.  For which one is the 
standard deviation of the numbers biggest?  smallest?

\hspace{5pt} (i)\hphantom{ii}  0, 20, 40, 50, 60, 80, 100\medskip

\hspace{5pt} (ii)\hphantom{i}  0, 48, 49, 50, 51, 52, 100\medskip

\hspace{5pt} (iii)             0, 1, 2, 50, 98, 99, 100
\bigskip


2. Guess which of the following two lists has the larger standard deviation.  Check your guess by 
computing the sd for both lists.

\hspace{5pt} (i)\hphantom{i} 9, 9, 10, 10, 12\medskip

\hspace{5pt} (ii)            7, 8, 10, 11, 13
\vfill
\eject

3. True or false:  the mean is always greater than the standard deviation?  
Find an example to justify your answer.
\vspace{1.2in}

4. True or false and explain:  SD and sd can't be negative.
\vspace{1.2in}

5. What are the SD and sd of the list 5, 5, 5, 5, 5?
\vspace{1.3in}

6. Find the sd for each list.  Explain how lists b)--d) are related to list a) and what impact
that has on the sd.

\hspace{10pt} a) 1, 3, 4, 4, 5, 7
\vspace{1.3in}

\hspace{10pt} b) 11, 13, 14, 14, 15, 17
\vspace{1.3in}

\hspace{10pt} c) 2, 6, 8, 8, 10, 14
\vspace{1.3in}

\hspace{10pt} d) $-3$, $-9$, $-12$, $-12$, $-15$, $-7$
\vfill
\eject
\end{document}
